\chapter{Grenzwerte}

\begin{definition}[Grenzwert]
  Sei $f: D \to \R$ eine Funktion mit $D \subseteq \R$ und $a$ ein Häufungspunkt von $D$. Dann heißt $L$ der Grenzwert von $f$ an der Stelle $a$, wenn für alle $\varepsilon > 0$ ein $\delta > 0$ existiert, so dass für alle $x \in D$ mit $0 < |x - a| < \delta$ gilt: $|f(x) - L| < \varepsilon$. Wir schreiben dann:
  \[
    \lim_{x \to a} f(x) = L.
  \]
\end{definition}

\begin{example}
  Betrachten wir die Funktion $f(x) = 2x + 3$. Wir wollen den Grenzwert von $f$ an der Stelle $a = 1$ bestimmen. Wir vermuten, dass $\lim_{x \to 1} f(x) = 5$ ist, da $f(1) = 2 \cdot 1 + 3 = 5$.

  Um dies zu zeigen, nehmen wir ein beliebiges $\varepsilon > 0$ und müssen ein $\delta > 0$ finden, so dass für alle $x$ mit $0 < |x - 1| < \delta$ gilt: $|f(x) - 5| < \varepsilon$.

  Berechnen wir $|f(x) - 5|$:
  \[
    |f(x) - 5| = |(2x + 3) - 5| = |2x - 2| = 2|x - 1|.
  \]

  Um sicherzustellen, dass $|f(x) - 5| < \varepsilon$, setzen wir $2|x - 1| < \varepsilon$, was äquivalent zu $|x - 1| < \frac{\varepsilon}{2}$ ist.

  Daher können wir $\delta = \frac{\varepsilon}{2}$ wählen. Dann gilt für alle $x$ mit $0 < |x - 1| < \delta$:
  \[
    |f(x) - 5| = 2|x - 1| < \varepsilon.
  \]

  Somit haben wir gezeigt, dass $\lim_{x \to 1} f(x) = 5$.
\end{example}

\begin{exercise}
  Berechnen Sie den Grenzwert $\lim_{x \to 0} \frac{\sin(x)}{x}$.
\end{exercise}

\begin{proposition}
  Sei $f: D \to \R$ eine Funktion mit $D \subseteq \R$ und $a$ ein Häufungspunkt von $D$. Wenn $\lim_{x \to a} f(x) = L$ existiert, dann ist dieser Grenzwert eindeutig.
\end{proposition}

\begin{proof}
  Angenommen, es existieren zwei verschiedene Grenzwerte $L_1$ und $L_2$ von $f$ an der Stelle $a$. Das bedeutet, dass für alle $\varepsilon > 0$ ein $\delta_1 > 0$ existiert, so dass für alle $x \in D$ mit $0 < |x - a| < \delta_1$ gilt: $|f(x) - L_1| < \varepsilon$. Ebenso existiert für alle $\varepsilon > 0$ ein $\delta_2 > 0$, so dass für alle $x \in D$ mit $0 < |x - a| < \delta_2$ gilt: $|f(x) - L_2| < \varepsilon$.

  Wählen wir nun $\varepsilon = \frac{|L_1 - L_2|}{3}$. Dann existieren $\delta_1 > 0$ und $\delta_2 > 0$, so dass für alle $x$ mit $0 < |x - a| < \delta_1$ gilt: $|f(x) - L_1| < \varepsilon$, und für alle $x$ mit $0 < |x - a| < \delta_2$ gilt: $|f(x) - L_2| < \varepsilon$.
\end{proof}

\begin{corollary}
  Wenn $f: D \to \R$ eine Funktion mit $D \subseteq \R$ und $a$ ein Häufungspunkt von $D$ ist, und wenn $\lim_{x \to a} f(x) = L$ existiert, dann gilt auch $\lim_{x \to a} |f(x)| = |L|$.
\end{corollary}

\begin{proof}
  Da $\lim_{x \to a} f(x) = L$ existiert, für alle $\varepsilon > 0$ ein $\delta > 0$ existiert, so dass für alle $x \in D$ mit $0 < |x - a| < \delta$ gilt: $|f(x) - L| < \varepsilon$. Wir wollen zeigen, dass $\lim_{x \to a} |f(x)| = |L|$.

  Berechnen wir $||f(x)| - |L||$:
  \[
    ||f(x)| - |L|| \leq |f(x) - L| < \varepsilon.
  \]

  Daher gilt für alle $x$ mit $0 < |x - a| < \delta$: $||f(x)| - |L|| < \varepsilon$. Somit haben wir gezeigt, dass $\lim_{x \to a} |f(x)| = |L|$.
\end{proof}

\begin{fact}
  Wenn $f: D \to \R$ eine Funktion mit $D \subseteq \R$ und $a$ ein Häufungspunkt von $D$ ist, und wenn $\lim_{x \to a} f(x) = L$ existiert, dann gilt auch $\lim_{x \to a} (f(x) + c) = L + c$ für jede Konstante $c \in \R$.
\end{fact}

\begin{theorem}
  Sei $f: D \to \R$ eine Funktion mit $D \subseteq \R$ und $a$ ein Häufungspunkt von $D$. Wenn $\lim_{x \to a} f(x) = L$ existiert, dann gilt auch $\lim_{x \to a} (c \cdot f(x)) = c \cdot L$ für jede Konstante $c \in \R$.
\end{theorem}

\begin{proof}
  Da $\lim_{x \to a} f(x) = L$ existiert, für alle $\varepsilon > 0$ ein $\delta > 0$ existiert, so dass für alle $x \in D$ mit $0 < |x - a| < \delta$ gilt: $|f(x) - L| < \varepsilon$. Wir wollen zeigen, dass $\lim_{x \to a} (c \cdot f(x)) = c \cdot L$.

  Berechnen wir $|c \cdot f(x) - c \cdot L|$:
  \[
    |c \cdot f(x) - c \cdot L| = |c| \cdot |f(x) - L| < |c| \cdot \varepsilon.
  \]

  Daher gilt für alle $x$ mit $0 < |x - a| < \delta$: $|c \cdot f(x) - c \cdot L| < |c| \cdot \varepsilon$. Somit haben wir gezeigt, dass $\lim_{x \to a} (c \cdot f(x)) = c \cdot L$.
\end{proof}